\documentclass[aspectratio=169]{beamer}

\usetheme{Madrid}
\usecolortheme{default}

\usepackage[utf8]{inputenc}
\usepackage[T1]{fontenc}
\usepackage[italian]{babel}
\usepackage{lmodern}
\usepackage{amsmath,amssymb,mathtools}
\usepackage{booktabs}
\usepackage{tikz}
\usetikzlibrary{positioning,arrows.meta,automata,calc}
\usepackage{listings}

\title{Input Driven Pushdown Automata ed Edit Distances}
\subtitle{Presentazione del lavoro di tesi}
\author{Luca Casnedi}
\institute{Universit\`a degli Studi di Milano\\Corso di Laurea in Informatica}
\date{Anno Accademico 2024--2025}

% --- DECORAZIONI GLOBALI ---
% Barra colorata sottile in alto
\setbeamertemplate{headline}{%
  \leavevmode
  \hbox{%
    \color{blue!60!black}\rule{\paperwidth}{1.5pt}%
  }%
}
% Linea sottile in basso
\setbeamertemplate{footline}{%
  \leavevmode
  \hbox{%
    \color{gray!60}\rule{\paperwidth}{0.5pt}%
  }%
}
% Margine slide leggermente aumentato per non "toccare" le barre
\addtobeamertemplate{frametitle}{}{\vspace{-0.5em}}

\begin{document}
\setbeamertemplate{navigation symbols}{}

\begin{frame}
  \titlepage
\end{frame}

\begin{frame}{Obiettivi della tesi}
\begin{itemize}
  \item Studiare il modello degli \emph{Input Driven Pushdown Automata} (IDPDA).
  \item Studiare il problema del calcolo della distanza di edit stringa-stringa e stringa-linguaggio.
  \item Evidenziare collegamenti tra automi input-driven e problema del calcolo delle distanze di edit.
  \item Evidenziare legami tra correction problem e questioni aperte in teoria della complessità.
\end{itemize}
\end{frame}

\begin{frame}{Motivazione}
Gli IDPDA estendono i DFA introducendo una pila con comportamento guidato dal simbolo letto.

\vspace{0.4em}
\begin{itemize}
  \item Adatti a strutture annidate e bilanciate.
  \item Riconoscono sottoclasse di linguaggi context-free.
  \item Operazione sulla pila determinata dal tipo di simbolo in input.
\end{itemize}
\end{frame}

\begin{frame}{Alfabeto input-driven}
Per un IDPDA, l'alfabeto si partiziona in tre insiemi disgiunti:
\[
\Sigma = \Sigma_{+1} \cup \Sigma_{-1} \cup \Sigma_0.
\]
\begin{itemize}
  \item $\Sigma_{+1}$: parentesi sinistre (operazione di \textit{push}).
  \item $\Sigma_{-1}$: parentesi destre (operazione di \textit{pop}).
  \item $\Sigma_0$: simboli neutri (nessuna operazione su pila).
\end{itemize}
\end{frame}

\begin{frame}{DIDPDA: definizione (idea)}
Un DIDPDA e' una settupla
\[
A=(\Sigma,Q,\Gamma,q_0,\perp,[\delta_a]_{a\in\Sigma},F)
\]
con transizioni determinate dalla categoria del simbolo in input.

\vspace{0.4em}
\begin{itemize}
  \item $\delta_{<}:Q\to Q\times\Gamma$ per $<\in\Sigma_{+1}$.
  \item $\delta_{>}:Q\times(\Gamma\cup\{\perp\})\to Q$ per $>\in\Sigma_{-1}$.
  \item $\delta_c:Q\to Q$ per $c\in\Sigma_0$.
\end{itemize}
\end{frame}

% Slide con diagramma a sinistra, divisorio nero centrale (fino in fondo) e liste a destra

\begin{frame}{Esempio di DIDPDA: diagramma}
\begin{columns}[T,onlytextwidth]
  \column{0.60\textwidth}
  \centering
  \begin{tikzpicture}[>=stealth, shorten >=1pt, auto, node distance=4cm, semithick, font=\scriptsize, scale=1.15, transform shape]
      \tikzstyle{every state}=[fill=white, draw=black, text=black, minimum size=12mm]
      \node[state, accepting] (q0) {$q_{0}$};
      \node[state, accepting, right=4cm of q0] (q1) {$q_{1}$};
      \node[state, below=3cm of $(q0)!0.5!(q1)$] (qerr) {$q_{\text{err}}$};
      % freccia iniziale dall'alto
      \draw[->] ([yshift=1cm]q0.north) -- (q0.north);
      % archi di transizione
      \path[->]
          (q0) edge[loop below] node[below left] {$<$/push $X$} ()
          (q0) edge[bend left=15] node[above] {$a$} (q1)
          (q0) edge node[left] {$>$/err} (qerr)
          (q1) edge[loop above] node[above] {$a$} ()
          (q1) edge[bend left=15] node[below] {$>$/pop $X$} (q0)
          (q1) edge node[right] {$>$/err} (qerr);
  \end{tikzpicture}%
  \hspace{2.5em}% Spazio aumentato tra diagramma e linea
  \column{0.015\textwidth}
  \begin{tikzpicture}[baseline={(current bounding box.center)}]
    \draw[black, thick] (0,0) -- (0,\dimexpr\textheight-2.5em\relax);
  \end{tikzpicture}
  \column{0.37\textwidth}
  \textbf{Stringhe accettate:}
  \begin{itemize}
    \item $<a>$
    \item $<<a>>$
    \item $a$
    \item $aa$
  \end{itemize}
  \vspace{0.5em}
  \textbf{Stringhe rifiutate:}
  \begin{itemize}
    \item $<<$
    \item $<a$
    \item $<>$
  \end{itemize}
\end{columns}
\vspace{0.5em}
\end{frame}

\begin{frame}{NIDPDA: estensione non deterministica}
Nel modello non deterministico, le transizioni restituiscono insiemi di stati possibili.

\vspace{0.4em}
\begin{itemize}
  \item Possibile insieme di stati iniziali $Q_0$.
  \item Per ogni simbolo, la computazione puo' biforcarsi.
  \item La stringa e' accettata se esiste almeno una computazione accettante.
\end{itemize}
\end{frame}

\begin{frame}{Grammatiche input-driven}
Le grammatiche input-driven usano lo stesso schema tripartito dell'alfabeto.

\vspace{0.4em}
Regole tipiche:
\[
A\rightarrow <B>C,\qquad A\rightarrow aC,\qquad A\rightarrow \epsilon.
\]

\begin{itemize}
  \item Origine storica: grammatiche di parentesi e grammatiche bilanciate.
  \item Vincolo centrale: corretta parentesizzazione.
\end{itemize}
\end{frame}

\begin{frame}{Distanza di edit tra stringhe}
Operazioni elementari: sostituzione, cancellazione, inserimento.

\vspace{0.4em}
\[
d(x,y)=\min\{\text{costo delle sequenze di edit che trasformano }x\text{ in }y\}.
\]

\begin{itemize}
  \item Costo complessivo = somma dei costi delle operazioni.
  \item Caso particolare: distanza di Hamming (solo sostituzioni).
\end{itemize}
\end{frame}

\begin{frame}{Distanza stringa-linguaggio e correction problem}
Per $L\subseteq\Sigma^*$ e $x\in\Sigma^*$:
\[
d(L,x)=\min\{d(y,x)\mid y\in L\}.
\]

\vspace{0.4em}
\textbf{Correction problem:}
\[
\text{trovare }y\in L\text{ tale che }d(y,x)=d(L,x).
\]

\begin{itemize}
  \item Non solo valutare una distanza, ma costruire una correzione ottima.
\end{itemize}
\end{frame}

\begin{frame}{Proprieta' utili per il calcolo}
Emergono due idee chiave:
\begin{itemize}
  \item \textbf{Decomposizione:}
  \[
  d(uv,wz)\le d(u,w)+d(v,z).
  \]
  \item \textbf{Limite superiore lineare:} esistono $\alpha,\beta$ tali che
  \[
  d(L,x)\le \alpha|x|+\beta.
  \]
\end{itemize}
Questo permette anche di limitare la lunghezza delle correzioni ottime.
\end{frame}

\begin{frame}{$k$-vicinato}
Definizione ricorsiva del vicinato di edit:
\[
E_0(w)=\{w\},\quad
E_{k+1}(w)=E_k(w)\cup\bigcup_{w'\in E_k(w)}\bigcup_{a\in\Sigma}(ins_a(w')\cup del_a(w')).
\]

Per un linguaggio:
\[
E_k(L)=\bigcup_{w\in L}E_k(w).
\]

\begin{itemize}
  \item Risultato: la famiglia dei linguaggi IDPDA e' chiusa rispetto al $k$-vicinato.
\end{itemize}
\end{frame}

\begin{frame}{Classi di complessita' considerate}
\begin{itemize}
  \item $P$: tempo polinomiale deterministico.
  \item $NP$: verifica polinomiale di certificati.
  \item $PSPACE$: spazio polinomiale.
  \item $AC^0$: circuiti di profondita' costante e taglia polinomiale.
\end{itemize}

Relazioni note:
\[
P\subseteq NP\subseteq PSPACE.
\]
\end{frame}

\begin{frame}{Teorema: collegamento con $P=NP$}
\begin{block}{Risultato}
Se $P=NP$, allora la distanza di edit di qualunque linguaggio in $NP$ e' calcolabile in tempo polinomiale.
\end{block}

Idea della prova:
\begin{itemize}
  \item si considera il problema decisionale $\Pi_L$: ``$d(L,x)\le k$?'';
  \item si indovina una correzione candidata di lunghezza limitata;
  \item si verifica appartenenza a $L$ e distanza in tempo polinomiale.
\end{itemize}
\end{frame}

\begin{frame}{Teorema: collegamento con $P=PSPACE$}
\begin{block}{Risultato}
Esiste un linguaggio context-sensitive $L$ tale che
\[
\text{edit-distance verso }L\in P\iff P=PSPACE.
\]
\end{block}

Interpretazione:
\begin{itemize}
  \item calcolare la distanza generalizza il test di appartenenza ($d(x,L)=0$);
  \item quindi, per opportuni linguaggi ``difficili'', la complessita' del problema riflette separazioni tra classi.
\end{itemize}
\end{frame}

\begin{frame}{Conclusioni}
\begin{itemize}
  \item Gli IDPDA offrono un compromesso potente tra espressivita' e struttura.
  \item La distanza di edit estende il riconoscimento a scenari rumorosi.
  \item I risultati di complessita' mostrano limiti profondi del problema.
  \item Aperture: algoritmi pratici di correzione, euristiche e varianti di costo.
\end{itemize}
\end{frame}

\begin{frame}{Riferimenti principali}
\footnotesize
\begin{thebibliography}{9}
\bibitem{okhotin-salomaa-2014}
A. Okhotin, K. Salomaa, \emph{Complexity of Input-Driven Pushdown Automata}, 2014.

\bibitem{pighizzini-2001}
G. Pighizzini, \emph{How Hard Is Computing the Edit Distance?}, 2001.

\bibitem{okhotin-salomaa-2018}
A. Okhotin, K. Salomaa, \emph{Edit distance neighbourhoods of input-driven pushdown automata}, 2018.
\end{thebibliography}
\end{frame}

\end{document}

